\chapter{国内外在该方向的研究现状及分析}


当前，联邦学习\cite{konevcny2015federated,mcmahan2017communication}依然是分布式计算领域中应用最为广泛的方法。然而，在分布式计算环境中，由于边缘设备的计算性能差且模型计算负担较重，联邦学习\cite{konevcny2015federated,mcmahan2017communication}面临着计算缓慢的问题。作为联邦学习与分割学习\cite{gupta2018distributed}的结合，分割-联邦学习针对这一问题提出了有效的解决方案。分割-联邦学习通过将模型分割并实施分布式计算，减轻了各节点的计算负担，显著降低了对边缘设备计算能力的需求。目前国内外有许多关于数据对分布式学习影响的研究，但关于分割-联邦学习在这一领域的研究仍处于起步阶段。
\section{分割-联邦学习（SplitFed）}
在分割-联邦学习提出之前，联邦学习\cite{konevcny2015federated,mcmahan2017communication}和分割学习\cite{gupta2018distributed}是最常用的分布式学习方法。联邦学习允许多个客户在不共享本地数据的情况下共同训练模型。参与联邦学习的客户在本地用私有的数据训练模型后仅共享本地模型的参数用来更新服务器中的全局模型，而不用共享其拥有的训练数据。这种训练模式可以有效保护用户隐私，特别适合于数据隐私和安全性要求高的场景。分割学习则是一种将深度学习模型分成两部分进行训练的方法，其中一部分在客户端进行计算，而另一部分在服务器上进行计算。这种方式可以有效减轻客户的计算负担，尤其是在训练数据量大且客户端设备能力有限的情况下。

联邦学习和分割学习都只需要一个服务器参与训练。在联邦学习中，服务器的任务是聚合不同用户的本地模型来得到全局模型，被称为聚合服务器；在分割学习中，服务器的任务是训练后半部分模型，被称为训练服务器。Chandra Thapa等人将分割学习中的训练服务器添加到联邦学习的系统中，提出了分割-联邦学习框架\cite{thapa2022splitfed}。该框架将模型分为前后两个部分，前半部分模型在客户本地和聚合服务器中按照联邦学习的方法进行训练，后半部分模型在计算服务器中按照分割学习的方法进行训练。

尽管 Chandra Thapa 等人提出的分割-联邦学习框架结合了联邦学习和分割学习的优点，但在准确率提升方面并未显著超越传统的联邦学习方法。这表明，分割-联邦学习在利用分布式计算资源进行模型训练时，未能有效改善模型性能。此外，分割-联邦学习同样面临着非独立同分布数据和标签噪声的影响，这些问题在分布式学习场景中普遍存在。具体来说，非独立同分布数据会导致模型在不同客户端间存在较大的分布差异，影响模型的泛化能力，而标签噪声则可能导致误导性更新，进一步降低了模型的收敛性和准确率。因此，如何在分割-联邦学习框架下有效处理这两种问题，以提升模型的整体性能，仍然是一个亟待解决的挑战。
\section{非独立同分布（non-IID）数据对分布式学习的影响}
非独立同分布（non-IID）数据在分布式学习中会显著影响模型的性能，尤其是在多个参与方的数据特征和分布差异较大的情况下\cite{zhu2021federated,amiri2022impact}。
现实世界中，非独立同分布数据的来源广泛，例如，社交媒体中的数据常常因个体差异而表现出显著的分布偏差。在医疗、金融等领域，由于不同机构的数据采集标准和样本选择的差异，也容易导致数据的非独立同分布特性。


非独立同分布数据会影响模型的收敛。Xiang Li等人证明了非独立同分布数据下联邦学习在光滑和强凸问题中的收敛时间复杂度\cite{li2019convergence}。
Sai Praneeth Karimireddy等人证明了非独立同分布数据会在联邦学习中引起客户端漂移（Client-Drift）从而导致
模型收敛速度慢和收敛过程不稳定\cite{karimireddy2020scaffold}。

目前有一些工作通过修改聚合函数构建了对非独立同分布数据具有鲁棒性的分布式学习算法。Tian Li在FedAvg的基础上引入了一个正则化项，提出了FedProx\cite{li2020federated}。
FedProx通过限制本地模型偏离全局模型的程度，帮助应对非独立同分布数据和异构设备。Sashank J. Reddi等人提出了FedAdam\cite{reddi2020adaptive}，这是一种
基于Adam优化器的联邦学习算法。其使用自适应优化方法，对全局更新使用一阶和二阶动量信息，以加速收敛，对于带有非独立同分布数据的复杂任务有一定的提升效果。
Di Wu等人提出了FedAdapt\cite{wu2022fedadapt}，它会动态调整聚合策略以适应不同客户端的模型更新步幅，以更好地处理高度非独立同分布的数据。


基于分割-联邦学习的衍生算法也能够在非独立同分布数据下有很好的表现。Yang Xu等人提出了SemiSFL\cite{xu2023semisfl}，该框架可以在非独立同分布数据和未标注数据共同存在的情况下实现模型的收敛。

尽管已有研究证明了非独立同分布数据对模型收敛性的影响，并提出了一些有效的对策，但针对分割-联邦学习框架在非独立同分布数据下的收敛性，目前尚缺乏系统性的理论证明。此外，尽管存在一些基于分割-联邦学习的衍生算法，如 SemiSFL，能够在非独立同分布数据和未标注数据共存的情况下实现模型的收敛，但仍未有研究提出能够显著提升分割-联邦学习在非独立同分布环境下准确率的算法。这表明，分割-联邦学习在应对非独立同分布数据挑战时，尚未充分挖掘其潜力，亟需进一步的探索和改进，以增强其在复杂数据环境中的表现。
\section{标签噪声对分布式学习的影响}
标签噪声会对分布式学习的模型的准确率产生显著影响，因为错误的标签
会误导模型的学习过程，使其在训练中偏离真实模式的识别\cite{liang2023fednoisy,jiang2024fnbench}。
例如，在分类任务中，含有噪声标签的数据会导致模型在学习
过程中混淆类别特征，从而降低对真实标签的区分能力，进而影响模
型的泛化性能。此外，标签噪声可能引发模型过拟合噪声样本，特别是
在数据量有限的情况下，模型会倾向于“记住”带有错误标签的数据，
导致在测试集上的准确率下降。目前，在深度学习中有许多减轻标签噪声
对模型影响的算法，目前很多方法已经被应用到分布式学习中。

Jingyi Xu等人提出了FedCorr\cite{xu2022fedcorr}。FedCorr基于
联邦学习的FedAvg聚合方法进行设计，能够对标签噪声
进行检测和修正。除此之外，FedCorr对不均匀的数据也具有较强的鲁棒性，
在复杂数据环境中的准确率超过同类算法20\%以上。

最近，也有很多对标签噪声具有鲁棒性的分布式算法被提出，这些新算法
均可以在擅长的场景下获得比FedCorr更高的准确率。其中包括Nannan Wu
提出的FedNoRo\cite{wu2023fednoro}，Jichang Li提出的FedDiv\cite{li2024feddiv}，Xinyuan Ji提出的FedFixer
\cite{ji2024fedfixer}
和Xiaochen Zhou等人提出的FedLNL\cite{zhou2024federated}。
等。 

尽管 Jingyi Xu 等人提出的 FedCorr 在标签噪声检测和修正方面表现出色，并在复杂数据环境中显著提升了准确率，但目前关于标签噪声对分布式学习收敛性影响的理论证明仍然十分稀缺。尤其是在分割-联邦学习框架下，相关研究几乎不存在。这表明，针对标签噪声对分割-联邦学习收敛性影响的系统性分析仍然是一个未被充分探讨的领域。此外，尽管近年来有多种基于其他框架的算法（如 FedNoRo、FedDiv 和 FedFixer）在噪声标签环境中表现良好，但至今尚未有研究提出基于分割-联邦学习的算法来提高其对标签噪声的鲁棒性。这一缺乏不仅限制了分割-联邦学习在复杂数据环境中的应用潜力，也使得提升其性能的机会亟需更多的研究和探索。




% \section{研究目的和意义}